\documentclass[11pt,a4paper]{article}
\usepackage{graphicx} 
\usepackage[left=2cm, right=2cm, top=2cm, bottom=2cm]{geometry}
\usepackage{amssymb}
\usepackage{amsmath}
\usepackage{setspace}
\setstretch{1.2}
\author{}
\date{}
\usepackage{fancyhdr}
\setlength{\headheight}{13.5pt}
\pagestyle{fancy}
\fancyhf{}
\fancyfoot[C]{-\thepage-}
\fancyhead[R]{}
\renewcommand{\headrulewidth}{0pt}
\renewcommand{\footrulewidth}{1pt}
\usepackage[colorlinks=true,linkcolor=blue,citecolor=red,urlcolor=green]{hyperref}
\usepackage[style=phys,backend=biber,biblabel=brackets]{biblatex}
\addbibresource{ref.bib}
\usepackage{bookmark}
\renewcommand{\thefootnote}{\fnsymbol{footnote}}
\begin{document}

\begin{center}
\begin{huge}
        \textbf{Synchrotron Radiation}
\end{huge} \\ \vspace{0.3cm}
\begin{small}
    - Rubya Akter Badhon    
\end{small}
\end{center}

\section*{Summary} EM radiation emitted by the charged particles when accelerated by force to move into a curved path at relativistic speed. 
So the e/m experiment that we have done in our lab, that does not count because the speed was not relativistic. But there was a glowing path. 
The light was from the gas, not from the accelerated electons. Synchrotron radiation is the dominent loss of energy in LEP collision experiment.
%--------------------------------------------------------------------------------------------------------------------------------------------------------

\section*{Setup}
The goal is to find the radiate power from an accelerated charge. 
This can be obtained from the poynting vector, which is the flow of energy per unit area per unit time or power per unit time. 
The poynting vector is,
\begin{align*}
    \mathbf{S} = \frac{1}{\mu_0}(\mathbf{E} \times \mathbf{B})
\end{align*}
And from this the power over the whole area becomes,
\begin{align*}
    P = \oint \mathbf{S} \cdot d\mathbf{a}
\end{align*}
$\mathbf{E}$ and $\mathbf{B}$ at time $t$ can be written as,
\begin{align*}
    & \mathbf{E} = -\nabla \phi - \frac{\partial \mathbf{A}}{\partial t} \\
    & \mathbf{B} = \nabla \times \mathbf{A}
\end{align*}
Here the fields are produced by a single accelerated elctron and these fields are called Li\`enard-Wiechert fields. 
Now I have to find $\phi$ and $\mathbf{A}$, which can be obtained from the waves equations that are,
\begin{align*}
    \Box^2 \phi &= -\frac{\rho}{\epsilon_0} \\
    \Box^2 \mathbf{A} &= -\mu_0 \mathbf{J}
\end{align*}
Here $\rho$ is the charge density and $\mathbf{J}$ is the current density. 
For the source terms of these equations, i have to use retarded time $t_{ret}$ to find the fields that they will create later at time $t$. 
\footnote{I will use bold letters for vectors.}
\newpage
%----------------------------------------------------------------------------------------------------------------------------------

\section*{Retarded Potentials and Retarded Time}
Here for synchroton radiation we have to deal with the retarded potential and retarded time. 
Cause, the radiation of present time that needs to be calculated is actually the result of the acceleration of electrons at retarded time. \\
Wave equations,
\begin{align}
    \Box^2 \phi &= -\frac{\rho}{\epsilon_0} \label{eq1} \\
    \Box^2 \mathbf{A} &= -\mu_0 \mathbf{J} \label{eq2}
\end{align}
Here $\Box^2 = (\nabla^2 - \frac{1}{c^2}\partial_t^2)$ is the d'alembertian and if $G(\mathbf{r},t;\mathbf{r'},t')$ is the green's function then,
\begin{align}
    &\Box^2 G(\mathbf{r},t;\mathbf{r'},t') = -\delta^3(\mathbf{r}-\mathbf{r'})\delta(t-t')\notag \\
    \implies & (\nabla^2 - \frac{1}{c^2}\partial_t^2)G(\mathbf{r},t;\mathbf{r'},t') = -\delta^3(\mathbf{r}-\mathbf{r'})\delta(t-t')
    \label{eq3}
\end{align}
As it is invariance in space time after translation, from equation (\ref{eq1}) and (\ref{eq3}),
\begin{align}
    \phi(r,t) = \frac{1}{\epsilon_0} \int G(\mathbf{r}-\mathbf{r'},t-t') \rho(\mathbf{r'},t')d^3r'dt' \label{eq4}
\end{align}
If $R\equiv (\mathbf{r}-\mathbf{r'})$ and $\tau \equiv (t-t')$ then using fourier's theorem,
\begin{align}
    G(\mathbf{R},\tau) = \int \frac{d\omega}{2\pi}\tilde{G}(\mathbf{R}, \omega) e^{-i\omega \tau} \label{eq5}
\end{align}
From equation (\ref{eq3}) and (\ref{eq5}) as translation happened as $\partial_t^2 \rightarrow -\omega^2$,
\begin{align}
     (\nabla^2 + \frac{\omega^2}{c^2})\tilde{G}(\mathbf{R}, \omega) = -\delta^3(\mathbf{R}) \label{eq6}
\end{align}
In this equation $\delta^3(\mathbf{R})$ is the source term (a point source at origin) and also everything else in equation (\ref{eq6}) is also rotationally symmetric. 
So $\mathbf{R} = |\mathbf{R}|$ and the solution of $\tilde{G}(R, \omega)$ can be written as,
\begin{align}
    \tilde{G}(R, \omega) = a\frac{e^{ikR}}{R} + b\frac{e^{-ikR}}{R} \label{eq7}
\end{align}
Using fourier's theorem again in equation (\ref{eq7}),
\begin{align}
    G(R, \tau) = \frac{a}{R} \delta(\tau - \frac{R}{c}) + \frac{b}{R} \delta(\tau + \frac{R}{c}) \label{eq8}
\end{align}
In this equation, there are two parts. now if we consider the causality we can have solution for one part, that is the retarded part. 
Due to causality, $G=0$ when $\tau<0$. 
 Now if $\tau = -\frac{R}{c}$ which is $<0$,
\begin{align}
     & G(R, -\frac{R}{c}) = \frac{a}{R}\delta(-\frac{R}{c}-\frac{R}{c}) + \frac{b}{R}\delta (-\frac{R}{c}+\frac{R}{c}) \notag \\
    \implies & G(R, -\frac{R}{c}) = \frac{b}{R} \tag{as $-2R/c \neq 0$ for $R>0$} \\
    \implies & \boxed{b = 0} \label{eq9}
\end{align}
From equation (\ref{eq7}) as $R\rightarrow 0$ for normalization,
\begin{align}
    \tilde{G}(R) = \frac{a}{R} + \frac{b}{R} \label{eq10}
\end{align}
Again near the origin as $R\rightarrow 0$, the exponential term of equation (\ref{eq7}) becomes 1.
\begin{align*}
     \nabla^2\tilde{G}(R) = -\delta^3(R)
\end{align*}
And the poisson's equation,
\begin{align} 
     & \nabla^2(\frac{1}{4\pi R}) = -\delta^3(R) \notag \\
     \implies & \tilde{G}(R) = \frac{1}{4\pi R} \notag \\
     \implies & a + b = \frac{1}{4\pi} \tag{from equation (\ref{eq10})} \\
     \therefore & \boxed{a = \frac{1}{4\pi}}
\end{align}
So the retarded Green's function can be obtain from equation (\ref{eq8}) as,
\begin{align}
    G_{ret}(R, \tau) = \frac{1}{4\pi R} \delta (\tau - \frac{R}{c}) \label{eq12}
\end{align}
and the scalar potential from equation (\ref{eq4}),
\begin{align}
    \phi_{ret}(r,t) = \frac{1}{4\pi \epsilon_0} \int \frac{\delta (t - t' -\frac{R}{c})}{R} \rho (r',t') d^3r'dt' \label{eq13}
\end{align}
From the delta function of equation (\ref{eq13}) the retarded time can be obtained as $t' = t_{ret}$.
\begin{align}
    \therefore \boxed{t_{ret} = t - \frac{R}{c}} \label{eq14}
\end{align}
Here $R = |\mathbf{r} - \mathbf{r'}|$. substituting them in equation (\ref{eq13}),
\begin{align}
    \therefore \boxed{\phi_{ret}(r,t) = \frac{1}{4\pi \epsilon_0} \int \frac{\rho (r',t_{ret})}{|\mathbf{r} - \mathbf{r'}|} d^3r'} \label{eq15}
\end{align}
This is the retarded scalar potential. Same goes for the retarded vector potential and from equation (\ref{eq2}) it can be written as,
\begin{align}
    \therefore \boxed{\mathbf{A}_{ret}(r,t) = \frac{\mu_0}{4\pi} \int \frac{\mathbf{J}(r',t_{ret})}{|\mathbf{r} - \mathbf{r'}|} d^3r'} \label{eq16}
\end{align}
\newpage
%---------------------------------------------------------------------------------------------------------------------------------

\section*{Li\`enard-Wiechert Fields}
For a point charge $q$ moving along $r_0(t)$ whose velocity is $\dot{r_0}(t') = \mathbf{v(t')}$ the density terms are,
\begin{align}
    \rho(r', t') &= q\delta^3(r' - r_0(t')) \label{eq17} \\
    \mathbf{J}(r', t') &= q\mathbf{v}(t')\delta^3(r' - r_0(t')) \label{eq18}
\end{align}
So if we calculate $\rho$, we can easily get $\mathbf{J}$. Substituting equation (\ref{eq17}) in equation (\ref{eq15}),
\begin{align}
    \phi_{ret}(r,t) &= \frac{1}{4\pi \epsilon_0} \int \frac{\delta (t - t' -\frac{|\mathbf{r} - \mathbf{r'|}}{c})}{|\mathbf{r} - \mathbf{r'}|}  q\delta^3(r' - r_0(t'))d^3r'dt' \notag
\end{align}
If $r' = r_0(t')$ then, 
\begin{align}
        \phi_{ret}(r,t) &= \frac{q}{4\pi \epsilon_0} \int \frac{\delta (t - t' -\frac{|\mathbf{r} - \mathbf{r_0}(t')|}{c})}{|\mathbf{r} - \mathbf{r_0}(t')|} dt' \notag \\
        &=  \frac{q}{4\pi \epsilon_0} \int \frac{\delta (t - t' -\frac{|\mathbf{R}|}{c})}{|\mathbf{R}|} dt' \label{eq19}
\end{align}
Here, the delta function contains the term which depends on $t'$. So, the function inside $\delta$ can be written as $f(t')$ from which
$f'(t')$ can be obtained.
\begin{align}
    f(t') &= t - t' - \frac{|\mathbf{r}-\mathbf{r_0}(t')|}{c} \notag \\
    f'(t') & = -1 - \frac{1}{c} \frac{d}{dt'}|\mathbf{r} - \mathbf{r_0}(t')| \notag \\
           & = -1 - \frac{1}{c} \frac{d}{dt'}|\mathbf{R}| \label{eq20}
\end{align}
and we know that,
\begin{align}
    \delta(f(t')) = \frac{\delta(t' - t_0)}{|f'(t_0)|} \label{eq21}
\end{align}
Again,
\begin{align}
    & |\mathbf{R}|^2 = \mathbf{R}\cdot \mathbf{R} \notag \\
    \therefore & \frac{d}{dt'}|\mathbf{R}| = \mathbf{R}\cdot \frac{\dot{\mathbf{R}}}{|\mathbf{R}|} \label{eq22} 
\end{align}
Therefore,
\begin{align}
    f'(t') = -1 + \frac{1}{c}\mathbf{R} \cdot \frac{\mathbf{v}}{|\mathbf{R}|} = -1 + \frac{\mathbf{v}}{c}\hat{R} \label{eq23}
\end{align}
At retarded time,
\begin{align}
    |f'(t_{ret})| = 1-\frac{\mathbf{v}}{c}.\hat{R}|_{t_{ret}} \label{eq24}
\end{align}
So, the scalar potential becomes,
\begin{align}
    \phi_{ret}(r,t) &= \frac{q}{4\pi \epsilon_0} \int \frac{1}{|\mathbf{r} - \mathbf{r_0}(t')|}\cdot \frac{\delta(t'-t_{ret})}{|f'(t_{ret})|} dt' \notag \\
    \therefore \phi_{ret}(r,t) &= \frac{1}{4\pi \epsilon_0}\cdot \frac{q}{R|f'(t_{ret})|} \label{eq25}
\end{align}
Here equation (\ref{eq25}) is obtained for $t' = t_{ret}$. Using similar method from equation (\ref{eq18}) the vector potential can also be obtained as,
\begin{align}
    \mathbf{A}_{ret}(r,t) &= \frac{\mu_0}{4\pi }\cdot \frac{q\mathbf{v}}{R|f'(t_{ret})|} \label{eq26}
\end{align}
Here equations (\ref{eq25}) and (\ref{eq26}) are called the Li\`enard-Wiechert potentials. 
Now that the potentials are known, the fields can also be calculated as,
\begin{align}                                                                                                                    
     & \mathbf{E} = -\nabla \phi - \frac{\partial \mathbf{A}}{\partial t}\label{eq27} \\                                                       
     & \mathbf{B} = \nabla \times \mathbf{A}\label{eq28}                                                                                       
\end{align} 
Another problem here is that the Li\`enard-Wiechert potentials of equation (\ref{eq25}) and (\ref{eq26}) are function of retarded time. 
But equation (\ref{eq27}) and (\ref{eq28}) depends on t. Have to get rid of $t_{ret}$...\\
Derivatives of $\phi$ can be written as,
\begin{align}
    \frac{\partial \phi}{\partial t} &= \frac{\partial \phi}{\partial t_{ret}} \cdot \frac{\partial t_{ret}}{\partial t}\label{eq29} \\
    \nabla \phi &= \frac{\partial \phi}{\partial t_{ret}}\cdot \nabla t_{ret} + \nabla \phi|_{t_{ret}} \label{eq30}  
\end{align}
These equations can be used to calculated electric and magnetic field without dealing with $\mathbf{A}$ as $\mathbf{A} = \frac{\mathbf{v}}{c^2}\phi$. \\
\begin{minipage}[c]{0.70\textwidth}
\begin{align}
    & t_{ret} = t-\frac{R(t_{ret})}{c} \notag \\
    & \frac{\partial t_{ret}}{\partial t} = 1-\frac{1}{c} \frac{\partial R}{\partial t_{ret}} \cdot \frac{\partial t_{ret}}{\partial t} \notag \\
    \implies & (1 - \frac{1}{c} \mathbf{v}\cdot \hat{R}) \frac{\partial t_{ret}}{\partial t} = 1 \notag 
\end{align}
\end{minipage}
\hspace{0.5em}\vrule height 3em \hspace{0.02em} \vrule height 2em \hspace{1em}
\begin{minipage}[c]{0.18\textwidth}
\begin{align*}
    \mathbf{R} &= \mathbf{r}-\mathbf{r_0}(t_{ret}) \\
    \frac{d\mathbf{R}}{dt_{ret}} &= -\mathbf{v}\cdot \hat{R}
\end{align*}
\end{minipage} 

\begin{align}
    \therefore \boxed{\frac{\partial t_{ret}}{\partial t} = \frac{1}{(1 - \frac{1}{c} \mathbf{v}\cdot \hat{R})} =  \frac{1}{|f'(t_{ret})|}} \label{eq31}
\end{align}
Again,
\begin{align}
    & \nabla t_{ret} = -\frac{1}{c} \nabla R (t_{ret}) \notag \\
    \implies & \nabla t_{ret} = -\frac{1}{c} [\nabla R +\frac{d\mathbf{R}}{dt_{ret}}\nabla t_{ret}] \notag \\
    \implies & \nabla t_{ret} = -\frac{1}{c} [\nabla R +(-\mathbf{v}\cdot \hat{R})\nabla t_{ret}] \notag \\
    \implies & (1 - \frac{1}{c} \mathbf{v}\cdot \hat{R}) \nabla t_{ret} = -\frac{1}{c} \hat{R} \notag 
\end{align}
\begin{align}
    \therefore \boxed{\nabla t_{ret} = \frac{-\frac{1}{c} \hat{R}}{(1 - \frac{1}{c} \mathbf{v}\cdot \hat{R})} = -\frac{\hat{R}}{c|f'(t_{ret})|}} \label{32}
\end{align}
From equation (\ref{eq25}),
\begin{align}
    \frac{\partial \phi}{\partial t_{ret}} &= \frac{q}{4\pi \epsilon_0} \frac{\partial}{\partial t_{ret}}(\frac{1}{R|f'(t_{ret})|}) \notag \\
                                           &= \frac{q}{4\pi \epsilon_0} \frac{-\frac{\partial}{\partial t_{ret}}(R|f'(t_{ret})|)}{R^2|f'(t_{ret})|^2} \notag \\
                                           &= \frac{q}{4\pi \epsilon_0} \frac{-R\frac{\partial}{\partial t_{ret}}(|f'(t_{ret})|) - |f'(t_{ret})|\frac{\partial R}{\partial t_{ret}}}{R^2|f'(t_{ret})|^2} \notag \\
                                           &= \frac{q}{4\pi \epsilon_0} \frac{-R(-\dot{\hat{R}}\mathbf{v}/c - \hat{R}\mathbf{a}/c) - |f'(t_{ret})|(-\mathbf{v}\cdot \hat{R})}{R^2|f'(t_{ret})|^2} \label{eq33} 
\end{align}
In equation (\ref{eq33}), $\dot{R} = -\hat{R}\cdot \mathbf{v}$ and the term $\dot{\hat{R}}$ can be obtained as,
\begin{align}
    \hat{R} &= \frac{d}{dt_{ret}}\left(\frac{\mathbf{r} - \mathbf{r_0}(t_{ret})}{R}\right) \notag \\
            &= \frac{-\mathbf{v}-\hat{R}\dot{R}}{R} \notag \\
            &= \frac{-\mathbf{v}-\hat{R}(\hat{R}.\mathbf{v})}{R} \label{eq34}
\end{align}
From equation (\ref{eq33}),
\begin{align}
    \frac{\partial \phi}{\partial t_{ret}} &= \frac{q}{4\pi \epsilon_0} \frac{[-\mathbf{v} + \hat{R}(\hat{R}\cdot \mathbf{v})]\mathbf{v}/c - R\hat{R}\cdot \mathbf{a}/c - |f'(t_{ret})|(-\mathbf{v}\cdot \hat{R})}{R^2|f'(t_{ret})|^2} \notag \\
                                           &= \frac{q}{4\pi \epsilon_0} \frac{-v^2/c + (\hat{R}\cdot \mathbf{v})^2/c - R\hat{R}\cdot \mathbf{a}/c - (1-\mathbf{v}.\hat{R}/c)(-\mathbf{v}\cdot \hat{R})}{R^2|f'(t_{ret})|^2} \notag \\
                                           &= \frac{q}{4\pi \epsilon_0} \frac{-v^2/c + R\hat{R}\cdot \mathbf{a}/c + \mathbf{v}\cdot \hat{R}}{R^2|f'(t_{ret})|^2} \label{eq35}
\end{align}
Again,
\begin{align}
    \nabla \phi|_{t_{ret}} &= \frac{q}{4\pi \epsilon_0} \nabla \left( \frac{1}{R|f'(t_{ret})|} \right) \notag \\
                         &= -\frac{q}{4\pi \epsilon_0} \frac{\nabla (R|f'(t_{ret})|)}{R^2|f'(t_{ret})|^2} \notag \\
                         &= -\frac{q}{4\pi \epsilon_0} \frac{\nabla (R - \frac{\mathbf{R}\cdot \mathbf{v}}{c})}{R^2|f'(t_{ret})|^2} \notag \\
                         &= -\frac{q}{4\pi \epsilon_0} \frac{\hat{R} - \frac{\mathbf{v}}{c}}{R^2|f'(t_{ret})|^2} \label{eq36}
\end{align}
From $\phi$ the vector potential term is,
\begin{align}
\frac{\partial \mathbf{A}}{\partial t} = \frac{1}{c^2}\left[\frac{\mathbf{a}}{|f'(t_{ret})|}\phi + \frac{\mathbf{v}}{|f'(t_{ret})|}\frac{\partial \phi}{\partial t_{ret}}\right ] \label{eq37}
\end{align}
Let, $\mathbf{u} = c\hat{R} - \mathbf{v}$ as denoted in the book \cite{book}
\footnote{ Now I have to take all the pieces and work out the fields. 
But first I need to simplify things as the terms are making me hazy. 
I have chosen the notation for me, so that I can recognize them at a glance.... }
and,
\begin{align*}
    & K = \frac{q}{4\pi \epsilon_0} \\
    & X = -v^2/c + R\hat{R}\cdot \mathbf{a}/c + \mathbf{v}\cdot \hat{R} \\
    & Y = |f'(t_{ret})|
\end{align*}
Then using the above equations, the simplified electric field can be written as,
\begin{align}
    \mathbf{E} &= \frac{KX\hat{R}}{R^2cY^3} + K\frac{\mathbf{u}}{cR^2Y^2} - \frac{1}{c^2}\left[\frac{\mathbf{a}K}{RY^2} + \frac{\mathbf{v}KX}{R^2Y^3}\right] \notag \\
               &= \frac{K}{cRY^3}\left[\frac{X\hat{R}}{R} + \frac{\mathbf{u}Y}{R} - \frac{\mathbf{a}Y}{c} - \frac{\mathbf{v}X}{cR}\right] \notag \\
               &= \frac{K}{c^2R^2Y^3}\left[ X\hat{R}c + \mathbf{u}Yc -\mathbf{a}RY - \mathbf{v}X \right] \notag \\
               &= \frac{K}{c^2R^2Y^3} \left[ X(\hat{R}c - \mathbf{v}) + Y \mathbf{u}c - \mathbf{a}RY \right] \notag \\
               &= \frac{K}{c^2R^2Y^3} \left[ X\mathbf{u} + Y\mathbf{u}c - \mathbf{a}RY \right] \notag \\
               &= \frac{K}{c^2R^2Y^3} \left[ (X + Yc)\mathbf{u} - \mathbf{a}RY \right] \label{eq38}
\end{align}
Here,
\begin{align*}
    X+Yc &= -\frac{v^2}{c} + \frac{R\hat{R}\cdot \mathbf{a}}{c} + \mathbf{v}\cdot \hat{R} + c - \mahtbf{v} \cdot \hat{R} \notag \\
         &= \frac{c^2-v^2}{c} + \frac{R}{c}\hat{R}\cdot \mathbf{a}
\end{align*}
Therefore,
\begin{align}
    \therefore \mathbf{E} &= \frac{K}{c^2R^2Y^3} \left[ (\frac{c^2 - v^2}{c} + \frac{R}{c}\hat{R}\cdot \mathbf{a} - \frac{\hat{R}\cdot \mathbf{u}}{c}) \mathbf{u} - \mathbf{a}R \frac{\hat{R}\cdot \mathbf{u}}{c} \right] \notag \\
                          &= \frac{K}{c^3R^2Y^3}\left[ (c^2-v^2)\mathbf{u} + R\left( (\hat{R}\cdot \mathbf{a})\mathbf{u} - (\hat{R}\cdot \mathbf{u})\mathbf{a} \right) \right] \notag \\
                          &= \frac{K}{c^3R^2Y^3}\left[ (c^2-v^2)\mathbf{u} + R \hat{R} \times (\mathbf{u}\times \mathbf{a})\right] \label{eq39}
\end{align}
In equation (\ref{eq39}). the radiated field term is the term that contains the acceleration $\mathbf{a}$. So,
\begin{align}
\therefore \boxed{\mathbf{E_{rad}} = \frac{K}{c^3RY^3} \hat{R}\times (\mathbf{u} \times \mathbf{a})} \label{eq40} 
\end{align}
\newpage
Here,

\begin{minipage}[c]{0.70\textwidth}                                                                                               
\begin{align}
    \mathbf{B} &= \nabla \times \mathbf{A} \notag \\
               &= \frac{1}{c^2} \nabla \times (\phi \mathbf{v}) \notag \\
               &= \frac{1}{c^2} \left[\nabla \phi \times \mathbf{v} + \phi (\nabla \times \mathbf{v}) \right] \label{eq41}
\end{align}                                                                                                                       
\end{minipage}                                                                                                                    
\hspace{0.5em}\vrule height 3em \hspace{0.02em} \vrule height 2em \hspace{1em}                                                    
\begin{minipage}[c]{0.18\textwidth}                                                                                               
\begin{align*}                                                                                                                    
    \mathbf{A} &= \phi \frac{\mathbf{v}}{c^2} \\
    \phi &= \frac{K}{RY}
\end{align*}                                                                                                                      
\end{minipage} 
\\

Terms in equation (\ref{eq41}) can be expressed as,
\begin{align}
     \nabla \phi &= -\frac{KX\hat{R}}{R^2cY^3} - \frac{K\mathbf{u}}{cR^2Y^2} \\ \notag \\
     (\nabla \times \mathbf{v})_i &= \epsilon_{ijk} \partial_j \times \mathbf{v}_k  \notag \\
                                  &= \epsilon_{ijk} (\partial_j t_{ret}) \times \mathbf{a}_k  \notag \\
                                  &= \nabla t_{ret} \times \mathbf{a} \tag{$\nabla t_{ret} = -\frac{\hat{R}}{cY}$} \\
                                  &= -\frac{\hat{R}\times \mathbf{a}}{cY} \\ \notag \\
     \nabla \phi \times \mathbf{v} &= \left(-\frac{KX\hat{R}}{R^2cY^3} - \frac{K\mathbf{u}}{cR^2Y^2}\right) \times \mathbf{v} \notag \\
                                   &= -\frac{K}{R^2cY^3} \left[ X\hat{R} + Y\mathbf{u}\right] \times \mathbf{v} \notag \\
                                   &= -\frac{K}{R^2cY^3} \left[ X\hat{R} \times \mathbf{v} + Y\mathbf{u} \times \mathbf{v} \right] \notag \\
                                   &= -\frac{K}{R^2cY^3} \left[ X\hat{R} \times \mathbf{v} + Yc\hat{R} \times \mathbf{v} - \mathbf{v} \times \mathbf{v} \right] \tag{substituting $\mathbf{u}$} \\
                                   &= -\frac{K}{R^2cY^3} (X + cY) \hat{R} \times \mathbf{v} 
\end{align}
Now that all conponents are acquired, $\mathbf{B}$ can be calculated using equation (\ref{eq41}).
\begin{align}
    \mathbf{B} &= \frac{1}{c^2} \left[\frac{K}{R^2cY^3} (X + cY) \hat{R} \times \mathbf{v} + \frac{K}{RY}\left( -\frac{\hat{R} \times \mathbf{a}}{cY} \right) \right] \label{eq45}  
\end{align}
From equation (\ref{eq45}) taking the acceleration term apart,
\begin{align}
    \mathbf{B'}_{rad} &= \frac{1}{c^2}\left[\frac{K}{R^2cY^3} X \hat{R} \times \mathbf{v} + \frac{K}{RY}\left( -\frac{\hat{R} \times \mathbf{a}}{cY} \right) \right] \notag \\
                      &= -\left[\frac{K}{R^2c^3Y^3} (-\frac{v^2}{c} + \frac{R}{c}\hat{R}\cdot \mathbf{a} + \mathbf{v}\cdot \hat{R}) \hat{R} \times \mathbf{v} + \frac{K}{c^3RY^2} \hat{R} \times \mathbf{a} \right] \label{eq46}
\end{align}
\footnote{ At first I was using only the term with $\mathbf{a}$ of equation (\ref{eq45}) and did wrong calculation. 
But $X$ also contains acceleration. Need to be cautious there.}
Again, in equation (\ref{eq46}) there are terms without acceleration, but only accelerated terms are responsible for radiation. 
Omitting them the actual $\mathbf{B}_{rad}$ can be obtained as,
\newpage
\begin{align}
\mathbf{B}_{rad} &= -\frac{K}{Rc^4Y^3} \hat{R} \cdot \mathbf{a}(\hat{R}\times \mathbf{v}) - \hat{R} \times \mathbf{v} + \frac{K}{c^3RY^2} \hat{R} \times \mathbf{a} \right] \notag \\
                 &= -\frac{K}{Rc^4Y^3} [\hat{R} \cdot \mathbf{a}(\hat{R}\times \mathbf{v}) + cY(\hat{R} \times \mathbf{a})] \notag \\
                 &= \frac{K}{Rc^4Y^3} [\hat{R} \cdot \mathbf{a}(\hat{R}\times \mathbf{u}) + cY(\hat{R} \times \mathbf{a})] \tag{$\mathbf{u} = c\hat{R}-\mathbf{v} \Rightarrow \hat{R} \times \mathbf{u} = -\hat{R} \times \mathbf{v} $} \\
                 &= \frac{K}{Rc^4Y^3} [\hat{R} \cdot \mathbf{a}(\hat{R}\times \mathbf{u}) - (c - \hat{R}\cdot \mathbf{v})(\hat{R} \times \mathbf{a})] \notag \\
                 &= \frac{K}{Rc^4Y^3} [\hat{R} \cdot \mathbf{a}(\hat{R}\times \mathbf{u}) - \hat{R} \cdot \mathbf{u}(\hat{R}\times \mathbf{a})] \notag \\
                 &= \frac{1}{c} \hat{R} \times \frac{K}{Rc^3Y^3}\left[\mathbf{u}(\hat{R} \cdot \mathbf{a}) - \mathbf{a}(\hat{R} \cdot \mathbf{u}) \right] \notag \\
                 &= \frac{1}{c} \hat{R} \times \frac{K}{Rc^3Y^3} \hat{R} \times (\mathbf{u} \times \mathbf{a}) \notag \\
\therefore &\boxed{\mathbf{B}_{rad} = \frac{1}{c} \hat{R} \times \mathbf{E}_{rad}} \label{eq47}
\end{align}
Equation (\ref{eq40}) and (\ref{eq47}) are acceleration parts of the Li\`enard-Wiechert fields.
%-----------------------------------------------------------------------------------------------------------------------------------

\section*{Power}
Using the Li\`enard-Wiechert fields poynting vector for radiation can be calculated,
\begin{align}
    \mathbf{S}_{rad} &= \frac{1}{\mu_0} \mathbf{E}_{rad} \times \mathbf{B}_{rad} \notag \\
                     &= \frac{1}{\mu_0} \mathbf{E}_{rad} \times \frac{1}{c} \hat{R} \times \mathbf{E}_{rad} \notag \\
                     &= \frac{1}{\mu_0 c}[ |\mathbf{E}_{rad}|^2 \hat{R} - \mathbf{E}_{rad} (\hat{R} \cdot \mathbf{E}_{rad})] \tag{$\hat{R} \cdot \mathbf{E}_{rad} = 0$} \\
                     &= \frac{1}{\mu_0 c} |\mathbf{E}_{rad}|^2 \hat{R} \label{eq48}
\end{align}

\footnote{From this point on relativistic notations will be used...}
let $\beta = \frac{\mathbf{v}}{c}$ and Lorentz factor $\gamma = \frac{1}{\sqrt{1-\frac{v^2}{c^2}}}$.
Then $\mathbf{u} = c(\hat{R} - \beta)$, $\mathbf{a} = c\dot{\beta}$ and $Y=1-\hat{R}.\beta$
\begin{align}
    \therefore \mathbf{E}_{rad} &= \frac{K}{Rc^3Y^3} \hat{R} \times [c(\hat{R}-\beta) \times c\dot{\beta}] \notag \\
                                &= \frac{K}{RcY^3} \hat{R} \times [(\hat{R}-\beta) \times \dot{\beta}] \label{eq49}
\end{align}
Now power per solid angle per unit time $t$,
\begin{align}
    \therefore \frac{dP}{d\Omega}\biggl|_t &= R^2 \mathbf{S}\cdot \hat{R} \notag \\
                                           &= \frac{R^2}{\mu_0 c}\left( \frac{q}{4\pi \epsilon_0cRY^3} \right)^2[\hat{R} \times (\hat{R}-\beta)\times \dot{\beta}]^2 \hat{R}\cdot \hat{R} \notag \\
                                           &= \frac{q^2}{16\pi^2\mu_0 \epsilon_0^2c^3Y^6}[\hat{R}\times \{ (\hat{R} - \beta) \times \dot{\beta} \}]^2 \tag{$c = 1/\sqrt{\mu_0 \epsilon_0}$} \\
                                           &= \frac{q^2}{16\pi^2 \epsilon_0cY^6}[\hat{R}\times \{ (\hat{R} - \beta) \times \dot{\beta} \}]^2 \label{eq50}
\end{align}
\\
The energy that the charge loses is for per unit retarded time. Here,

\begin{minipage}[c]{0.60\textwidth}                                                                                               
 \begin{align}
     & t - t_{ret} = \frac{R}{c} \notag \\
     & \frac{dt}{dt_{ret}} -1 = \frac{\dot{R}}{c} \notag \\
     & \frac{dt}{dt_{ret}} = 1 - \hat{R}\cdot \frac{\mathbf{v}}{c} = Y \notag \\
     & dt = Ydt_{ret}
 \end{align}                                                                                                                       
\end{minipage}                                                                                                                    
\hspace{0.5em}\vrule height 3em \hspace{0.02em} \vrule height 2em \hspace{1em}                                                    
\begin{minipage}[c]{0.1\textwidth}                                                                                               
 \begin{align*}
     \dot{R} &= - \hat{R}\cdot \frac{\mathbf{v}}{c} \\
     Y &= 1 - \hat{R}\cdot \frac{\mathbf{v}}{c} \\
       &= 1 - \hat{R}\cdot \beta
 \end{align*}                                                                                                                      
\end{minipage}                                                                                                                    
\\
Then,
\begin{align}
    \therefore \frac{dP}{d\Omega}\biggl|_{t_{ret}} &= \frac{q^2}{16\pi^2 \epsilon_0cY^5}[\hat{R}\times \{ (\hat{R} - \beta) \times \dot{\beta} \}]^2 \label{eq52}
\end{align}
The equation contains a triple product which needs to be expended,
\begin{align}
    [\hat{R}\times \{ (\hat{R} - \beta) \times \dot{\beta} \}]^2 &= [(\hat{R} - \beta)(\hat{R} \cdot \dot{\beta}) - \dot{\beta}\{\hat{R}\cdot (\hat{R} - \beta)\}]^2 \notag \\
                                                                 &= [(\hat{R} - \beta)(\hat{R} \cdot \dot{\beta}) - Y\dot{\beta}]^2 \tag{$\because \hat{R}\cdot (\hat{R} - \beta) = Y$} \\
                                                                 &= (\hat{R} - \beta)^2(\hat{R} \cdot \dot{\beta})^2 + Y^2\dot{\beta}^2 - 2(\hat{R} - \beta)(\hat{R} \cdot \dot{\beta})Y\dot{\beta}  \notag \\
                                                                 &= (1 - 2\hat{R}\cdot \beta +\beta^2)(\hat{R} \cdot \dot{\beta})^2 + Y^2\dot{\beta}^2 - 2\hat{R}(\hat{R} \cdot \dot{\beta})Y\dot{\beta} + 2\beta (\hat{R} \cdot \dot{\beta})Y\dot{\beta} \notag \\
                                                                 &= (1 - 2(1-Y) +\beta^2)(\hat{R} \cdot \dot{\beta})^2 + Y^2\dot{\beta}^2 - 2Y(\hat{R} \cdot \dot{\beta})^2 + 2Y(\hat{R} \cdot \dot{\beta}) (\beta \cdot \dot{\beta}) \notag \\
                                                                 &= (2Y - \frac{1}{\gamma^2})(\hat{R} \cdot \dot{\beta})^2 + Y^2\dot{\beta}^2 - 2Y(\hat{R} \cdot \dot{\beta})^2 + 2Y(\hat{R} \cdot \dot{\beta}) (\beta \cdot \dot{\beta}) \notag \\
                                                                 &= -\frac{(\hat{R} \cdot \dot{\beta})^2}{\gamma^2} + Y^2\dot{\beta}^2 \label{eq53}
\end{align}
Equation (\ref{eq53}) is true for circular motion, where $\beta \perp \dot{\beta}$. 
In spherical coordinate system, let's take $\theta$ the angle between $\hat{R}$ and $\beta$.
Then $Y = 1 - \beta \cos{\theta}$ and $\hat{R}\cdot \dot{\beta} = \dot{\beta} \sin{\theta} \cos{\phi}$.
So, equation (\ref{eq52}) becomes,
\begin{align}
    \frac{dP}{d\Omega}\biggl|_{t_{ret}} &= \frac{q^2}{16\pi^2 \epsilon_0c} \left[ \frac{\dot{\beta^2}}{(1-\beta \cos{\theta})^3} - \frac{\dot{\beta}^2\sin^2{\theta}\cos^2{\phi}}{\gamma^2(1-\beta \cos{\theta})^5} \right] \notag \\
                                        &= \frac{q^2\dot{\beta}^2}{16\pi^2 \epsilon_0c} \left[ \frac{1}{(1-\beta x)^3} - \frac{(1-x^2)\cos^2{\phi}}{\gamma^2(1-\beta x)^5} \right] \notag
\end{align}
This equation is obtained by changing variable. 
Here, $x=\cos{\theta}$ and for that the limits become $\theta = -\pi \rightarrow x = -1$ and $\theta = \pi \rightarrow x = 1$.
And also $d\Omega = \sin{\theta}d\theta d\phi $
Then power over the whole area becomes,
\begin{align}                                                                                                                     
    \int \frac{dP}{d\Omega} d\Omega &= \frac{q^2\dot{\beta^2}}{16\pi^2 \epsilon_0c}\int_{0}^{2\pi} \int_{-1}^{1} \left[ \frac{1}{(1-\beta \cos{\theta})^3} - \frac{(1-x^2)\cos^2{\phi}}{\gamma^2(1-\beta x)^5} \right] dx d\phi \label{eq54}
\end{align}
There are two terms in this equation. 
I will integrate each term seperately. 
Integrating the 1st term,
\begin{align}
    \int_{0}^{2\pi} \int_{-1}^{1} \frac{1}{(1-\beta x)^3} dxd\phi &= 2\pi \frac{1}{2\beta}(1-\beta x)^{-2}\biggr|^1_{-1} \notag \\
                                                                  &= \frac{\pi}{\beta}\left( \frac{4\beta}{(1-\beta^2)^2} \right)
\end{align}
Now  for integration of the 2nd term, changing variable again,
\begin{align*}
    1 - \beta x = w &\rightarrow x = \frac{1-w}{\beta} \Rightarrow dx = \frac{1}{\beta}dw \\
    x = -1 &\rightarrow w = 1+\beta \\
    x = 1 &\rightarrow w= 1-\beta
\end{align*}
From equation (\ref{eq54}),
\begin{align}
    & \int_{0}^{2\pi} \cos^2{\phi}d\phi \int_{1+\beta}^{1-\beta} \frac{1}{\gamma^2 \beta} \frac{1-\frac{(1-w)^2}{\beta^2}}{w^5}dw \notag \\
    =& (\frac{1}{2}\phi + \frac{1}{4}\sin(2\phi))\biggr|_0^{2\pi} \int_{1+\beta}^{1-\beta} \frac{1}{\gamma^2 \beta} \frac{\frac{\beta^2-1 + 2w -w^2}{\beta^2}}{w^5} dw \notag \\
    =& \pi \int_{1+\beta}^{1-\beta} \frac{1}{\gamma^2 \beta} \frac{\frac{-\frac{1}{\gamma^2} + 2w -w^2}{\beta^2}}{w^5} dw \notag \\
    =& \pi \int_{1+\beta}^{1-\beta} \frac{1}{\gamma^2 \beta^3} \frac{-\frac{1}{\gamma^2} + 2w -w^2}{w^5} dw \notag \\
    =& \pi \int_{1+\beta}^{1-\beta} \frac{1}{\gamma^2 \beta^3}(-\frac{1}{\gamma^2}w^{-5} + 2w^{-4} - w^{-3}) dw \notag \\
    =& \frac{\pi}{\gamma^2 \beta}\left[ \frac{1}{4\gamma^2} (1+\beta)^{-4} -\frac{2}{3}(1+\beta)^{-3} +\frac{1}{2}(1+\beta)^{-2} - \frac{1}{4\gamma^2} (1-\beta)^{-4} +\frac{2}{3}(1-\beta)^{-3} -\frac{1}{2}(1-\beta)^{-2} \right] \notag
\end{align}
Again for simplification let, $ 1+\beta = x$ and $ 1-\beta = y$ then
\begin{align}
    xy = (1+\beta)(1-\beta) = 1 - \beta^2 = \frac{1}{\gamma^2} 
\end{align}
Then,
\begin{align}
    & \frac{\pi}{\gamma^2 \beta}\left[ \frac{1}{4\gamma^2} (1+\beta)^{-4} -\frac{2}{3}(1+\beta)^{-3} +\frac{1}{2}(1+\beta)^{-2} - \frac{1}{4\gamma^2} (1-\beta)^{-4} +\frac{2}{3}(1-\beta)^{-3} -\frac{1}{2}(1-\beta)^{-2} \right] \notag \\
    =& \frac{\pi}{\gamma^2 \beta}\left[ \frac{1}{4\gamma^2} (\frac{1}{x^4} - \frac{1}{y^4}) - \frac{2}{3} (\frac{1}{x^3} - \frac{1}{y^3}) + \frac{1}{2} (\frac{1}{x^2} - \frac{1}{y^2}) \right] \label{eq57}  
\end{align}
Let's calculate some terms of equation (\ref{eq57}) first,
\begin{align}
    \therefore & y^2-x^2 = (1 -2\beta +\beta^2 -1 -\beta^2 -2\beta) = -4\beta \\
    \therefore & y^4-c^4 = (y^2 - x^2)(y^2 + x^2) = -4\beta(2+2\beta^2) = -8\beta(1 +\beta^2) \\
    \therefore & y^3-x^3 = (y-x)(x^2 +y^2 +xy) = -2\beta (2 +2\beta^2 +1 -\beta^2) = -2\beta(3 +\beta^2)
\end{align}
Then equation (\ref{eq57}) becomes,
\begin{align}
    & \frac{\pi}{\gamma^2 \beta}\left[ \frac{1}{4}\gamma^6 (-8\beta) (1+\beta^2) - \frac{2}{3}\gamma^6 (-2\beta) (3+\beta^2) + \frac{1}{2}\gamma^4 (-4\beta) \right] \notag \\
    =& \pi \left[ -2\gamma^4 (1+\beta^2) + \frac{4}{3}\gamma^4 (3+\beta^2) -2\gamma^2 \right] \notag \\
    =& \pi \left[ -2\gamma^4 -2\gamma^4\beta^2 + 4\gamma^4+\frac{4}{3}\gamma^4 \beta^2 -2\gamma^2 \right] \notag
\end{align}
\begin{align}
    =& \pi \left[ -2\gamma^4 -2\gamma^4\beta^2 + 4\gamma^4+\frac{4}{3}\gamma^4 \beta^2 -2\gamma^2 \right] \notag \\
    =& \pi \left[ -2\gamma^4 -\frac{2}{3}\gamma^4\beta^2 -2\gamma^4(1-\beta^2) \right] \notag \\
    =& 2\gamma\pi \left[ 1-\frac{1}{3}\beta^2 -(1-\beta^2) \right] \notag \\
    =& \frac{4}{3}\pi\gamma^4
\end{align}
Summing everything up equation (\ref{eq54}) becomes,
\begin{align}
    \int \frac{dP}{d\Omega} d\Omega &= \frac{q^2\dot{\beta^2}}{16\pi^2 \epsilon_0c} (4\pi\gamma^4 - \frac{4}{3}\pi\gamma^4) \notag \\
                                    &= \frac{q^2\dot{\beta^2}}{4\pi^2 \epsilon_0c}(\frac{2}{3}\pi \gamma^4) \notag
\end{align}
Then the radiation power becomes, 
\begin{align}
    \therefore \boxed{P = \frac{q^2\dot{\beta}^2\gamma^4}{6\pi \epsilon_0c}} \label{eq62}
\end{align}
This is the \textbf{synchrotron radiation power:} the rate at which a relativistic charge loses energy to radiation. 
In this equation $\dot{\beta} = \frac{\mathbf{a}}{c}$, if the speed is not relativistic then $\gamma \rightarrow 1$ 
and for that equation (\ref{eq62}) becomes,
\begin{align}
    \therefore P = \frac{q^2a^2}{6\pi \epsilon_0c^3} \label{eq63} 
\end{align}
Equation (\ref{eq63}) is the Larmor equation, which is what expected for nonrelativistic case.  \\ \\
Again equation (\ref{eq62}) also depends on mass even though there is no direct mass term in it. 
It contain $\gamma^4$ and we know that, $\gamma = \frac{E}{mc^2}$. 
so, for fixed energy and radius, there is inversely proportional mass of the power of 4 in the power equation. 
That is why electron will radiate way more than proton ($\approx 1836^4$ times) at same energy.
In LEP experiment where electron-positron collision is used, the synchroton radiation is huge and it causes a lot of enery loss.
In LHC experiment the energy loss is much lower as there proton-proton collision is used.  
\nocite{*}
\printbibliography
\end{document}


 
